\documentclass{article}
\usepackage{amsmath,amssymb,amsthm}
\usepackage{algorithm,algorithmic}
\usepackage{hyperref}
\usepackage{enumitem}
\newtheorem{theorem}{Theorem}
\newtheorem{lemma}{Lemma}
\newtheorem{assumption}{Assumption}
\newcommand{\diam}{\operatorname{diam}}
\newcommand{\gap}{\operatorname{g}}
\begin{document}

%%%%%%%%%%%%%%%%%%%%%%%%%%%%%%%%%%%%%%%%%%%%%%%%%%%%%%%%%%%%
\section*{Frank–Wolfe (Conditional Gradient) Method}
%%%%%%%%%%%%%%%%%%%%%%%%%%%%%%%%%%%%%%%%%%%%%%%%%%%%%%%%%%%%

\subsection*{Mathematical Formulation}
Let $f:\mathbb{R}^n\rightarrow\mathbb{R}$ be a convex, differentiable function and let 
\[
\mathcal{C}\subset\mathbb{R}^n
\]
be a non‑empty compact convex set (the feasible domain).  
At iteration $k\ge 0$ we have a current point $x_k\in\mathcal{C}$.  
The algorithm queries the \emph{linear minimization oracle} (LMO)

\[
s_k \;=\; \arg\min_{s\in\mathcal{C}}\;\langle\nabla f(x_k),\,s\rangle .
\tag{1}
\]

A step size $\gamma_k\in[0,1]$ is chosen (either by a prescribed schedule or by exact line search) and the next iterate is a convex combination of $x_k$ and $s_k$:

\[
x_{k+1}\;=\;(1-\gamma_k)x_k+\gamma_k s_k .
\tag{2}
\]

The algorithm terminates when a stopping criterion (e.g. the Frank–Wolfe gap below a tolerance) is satisfied.

\subsection*{Pseudocode}
\begin{algorithm}[H]
\caption{Frank–Wolfe Method}
\begin{algorithmic}[1]
\REQUIRE Convex function $f$, feasible set $\mathcal{C}$, initial point $x_0\in\mathcal{C}$, 
step‑size rule $\{\gamma_k\}_{k\ge0}$, tolerance $\varepsilon>0$.
\FOR{$k=0,1,2,\dots$}
    \STATE Compute gradient $g_k \gets \nabla f(x_k)$.
    \STATE \textbf{LMO:} $s_k \gets \displaystyle\arg\min_{s\in\mathcal{C}}\langle g_k,s\rangle$.
    \STATE Compute the Frank–Wolfe gap 
          \[
          \gap_k \gets \langle g_k, x_k-s_k\rangle .
          \]
    \IF{$\gap_k \le \varepsilon$}
        \STATE \textbf{break}
    \ENDIF
    \STATE Choose step size $\gamma_k\in[0,1]$ (e.g. $\gamma_k = \frac{2}{k+2}$ or exact line search).
    \STATE Update $x_{k+1} \gets (1-\gamma_k)x_k + \gamma_k s_k$.
\ENDFOR
\RETURN $x_k$
\end{algorithmic}
\end{algorithm}

\subsection*{Dry Run Example}
Consider the one–dimensional problem  

\[
f(w) = (w-3)^2 ,\qquad \mathcal{C}=[0,5] .
\]

The gradient is $\nabla f(w)=2(w-3)$.  
We initialise $w_0=0$, fix a constant step size $\gamma=0.1$, and run three iterations.

\begin{center}
\begin{tabular}{c|c|c|c}
iteration $k$ & $w_k$ & $\nabla f(w_k)$ & $w_{k+1}=(1-\gamma)w_k+\gamma s_k$ \\ \hline
0 & $0$ & $-6$ & $s_0=\arg\min_{s\in[0,5]}(-6)s = 5$\\
  &   &      & $w_1 = 0.9\cdot0+0.1\cdot5 = 0.5$ \\ \hline
1 & $0.5$ & $-5$ & $s_1=\arg\min_{s\in[0,5]}(-5)s = 5$\\
  &   &      & $w_2 = 0.9\cdot0.5+0.1\cdot5 = 0.95$ \\ \hline
2 & $0.95$ & $-4.1$ & $s_2=\arg\min_{s\in[0,5]}(-4.1)s = 5$\\
  &   &      & $w_3 = 0.9\cdot0.95+0.1\cdot5 \approx 1.355$ \\ \hline
\end{tabular}
\end{center}

The corresponding objective values are  

\[
f(w_0)=9,\qquad f(w_1)=6.25,\qquad f(w_2)=4.2025,\qquad f(w_3)\approx 2.71 .
\]

The iterates clearly move toward the optimum $w^\star=3$ while remaining inside the feasible interval.

%%%%%%%%%%%%%%%%%%%%%%%%%%%%%%%%%%%%%%%%%%%%%%%%%%%%%%%%%%%%
\section*{Assumptions}
%%%%%%%%%%%%%%%%%%%%%%%%%%%%%%%%%%%%%%%%%%%%%%%%%%%%%%%%%%%%
\begin{assumption}[Convexity and Smoothness]\label{ass:convex-smooth}
The objective $f$ is convex and $L$‑smooth on $\mathcal{C}$, i.e. for all $x,y\in\mathcal{C}$,
\[
f(y)\le f(x)+\langle\nabla f(x),y-x\rangle+\frac{L}{2}\|y-x\|^2 .
\tag{3}
\]
\end{assumption}

\begin{assumption}[Bounded Feasible Set]\label{ass:bounded}
The feasible set $\mathcal{C}$ is compact with finite diameter
\[
D \;=\; \diam(\mathcal{C}) \;=\; \max_{x,y\in\mathcal{C}}\|x-y\| <\infty .
\tag{4}
\]
\end{assumption}

\begin{assumption}[Linear Minimization Oracle]\label{ass:lmo}
For any $x\in\mathcal{C}$ the LMO (1) can be solved exactly.
\end{assumption}

%%%%%%%%%%%%%%%%%%%%%%%%%%%%%%%%%%%%%%%%%%%%%%%%%%%%%%%%%%%%
\section*{Main Convergence Theorem}
%%%%%%%%%%%%%%%%%%%%%%%%%%%%%%%%%%%%%%%%%%%%%%%%%%%%%%%%%%%%
\begin{theorem}[Primal Gap $O(1/k)$]\label{thm:rate}
Let Assumptions \ref{ass:convex-smooth}–\ref{ass:lmo} hold and let the step size be
\[
\gamma_k \;=\; \frac{2}{k+2},\qquad k\ge0 .
\tag{5}
\]
Denote $x^\star\in\arg\min_{x\in\mathcal{C}}f(x)$.  
Then for all $k\ge0$,
\[
f(x_k)-f(x^\star) \;\le\; \frac{2L D^2}{k+2}.
\tag{6}
\]
Equivalently, the Frank–Wolfe (duality) gap satisfies
\[
\gap_k \;=\; \langle\nabla f(x_k),x_k-s_k\rangle \;\le\; \frac{4L D^2}{k+2}.
\tag{7}
\]
\end{theorem}

\subsection*{Theoretical Properties}
\begin{itemize}[leftmargin=*]
\item \textbf{Stability.} The iterates remain in $\mathcal{C}$ for all $k$ because (2) is a convex combination of two points in $\mathcal{C}$.
\item \textbf{Step–size sensitivity.} The canonical choice (5) yields the optimal $O(1/k)$ rate without line search. Any $\gamma_k\in[0,1]$ satisfying $\gamma_k\ge \frac{c}{k}$ for a constant $c>0$ also yields $O(1/k)$ convergence, albeit with a larger constant.
\item \textbf{Comparison to projected gradient.} Frank–Wolfe avoids Euclidean projections; the rate (6) matches the standard $O(1/k)$ rate of projected gradient for smooth convex problems, while each iteration may be cheaper when the LMO is inexpensive.
\end{itemize}

\subsection*{Discussion}
The bound (6) depends only on the smoothness constant $L$ and the geometry of $\mathcal{C}$ via its diameter $D$. No strong convexity is required. Moreover, the algorithm enjoys a natural \emph{sparsity} property: when $\mathcal{C}$ is a polytope, each $x_k$ is a convex combination of at most $k+1$ extreme points.

%%%%%%%%%%%%%%%%%%%%%%%%%%%%%%%%%%%%%%%%%%%%%%%%%%%%%%%%%%%%
\section*{Preliminaries (Definitions and Lemmas)}
%%%%%%%%%%%%%%%%%%%%%%%%%%%%%%%%%%%%%%%%%%%%%%%%%%%%%%%%%%%%
\begin{definition}[Curvature Constant]\label{def:curvature}
For a smooth convex $f$ on $\mathcal{C}$ define
\[
C_f \;:=\; \sup_{\substack{x,s\in\mathcal{C}\\ \gamma\in[0,1]}}
\frac{2}{\gamma^2}\Bigl(
f\bigl(x+\gamma (s-x)\bigr)-f(x)-\gamma\langle\nabla f(x),s-x\rangle
\Bigr).
\tag{8}
\]
\end{definition}
From smoothness (3) one obtains the simple bound $C_f\le L D^2$.

\begin{lemma}[Descent Lemma for Convex Smooth Functions]\label{lem:descent}
For any $x\in\mathcal{C}$, $s\in\mathcal{C}$ and $\gamma\in[0,1]$,
\[
f\bigl((1-\gamma)x+\gamma s\bigr)
\le f(x)+\gamma\langle\nabla f(x),s-x\rangle+\frac{L\gamma^2}{2}\|s-x\|^2 .
\tag{9}
\]
\end{lemma}
\begin{proof}
Apply the $L$‑smoothness inequality (3) with $y=(1-\gamma)x+\gamma s$ and use $\|y-x\|=\gamma\|s-x\|$.
\end{proof}

%%%%%%%%%%%%%%%%%%%%%%%%%%%%%%%%%%%%%%%%%%%%%%%%%%%%%%%%%%%%
\section*{Main Proof (Step‑by‑Step Derivation)}
%%%%%%%%%%%%%%%%%%%%%%%%%%%%%%%%%%%%%%%%%%%%%%%%%%%%%%%%%%%%
We prove Theorem \ref{thm:rate} by establishing a recurrence for the primal gap and then solving it.

\subsection*{Step 1 – Apply Lemma \ref{lem:descent}}
From (9) with $x=x_k$, $s=s_k$ and $\gamma=\gamma_k$ we have
\[
f(x_{k+1})
\le f(x_k)+\gamma_k\langle\nabla f(x_k),s_k-x_k\rangle
      +\frac{L\gamma_k^{2}}{2}\|s_k-x_k\|^{2}.
\tag{10}
\]
\underline{[Step 1]}

\subsection*{Step 2 – Use definition of the Frank–Wolfe gap}
Recall $\gap_k = \langle\nabla f(x_k),x_k-s_k\rangle = -\langle\nabla f(x_k),s_k-x_k\rangle$.
Thus (10) becomes
\[
f(x_{k+1}) \le f(x_k)-\gamma_k\gap_k
               +\frac{L\gamma_k^{2}}{2}\|s_k-x_k\|^{2}.
\tag{11}
\]
\underline{[Step 2]}

\subsection*{Step 3 – Bound the distance term}
Since $x_k,s_k\in\mathcal{C}$, by the diameter definition (4)
\[
\|s_k-x_k\|\le D .
\tag{12}
\]
Insert (12) into (11):
\[
f(x_{k+1}) \le f(x_k)-\gamma_k\gap_k+\frac{L D^{2}}{2}\gamma_k^{2}.
\tag{13}
\]
\underline{[Step 3]}

\subsection*{Step 4 – Relate gap to primal suboptimality}
Convexity of $f$ gives for any $x^\star\in\mathcal{C}$
\[
f(x_k)-f(x^\star) \le \langle\nabla f(x_k),x_k-x^\star\rangle .
\tag{14}
\]
Because $s_k$ minimizes $\langle\nabla f(x_k),s\rangle$ over $\mathcal{C}$,
\[
\langle\nabla f(x_k),s_k\rangle \le \langle\nabla f(x_k),x^\star\rangle .
\tag{15}
\]
Subtracting (15) from $\langle\nabla f(x_k),x_k\rangle$ yields
\[
\gap_k = \langle\nabla f(x_k),x_k-s_k\rangle
       \ge \langle\nabla f(x_k),x_k-x^\star\rangle .
\tag{16}
\]
Combining (14) and (16) we obtain
\[
\gap_k \;\ge\; f(x_k)-f(x^\star).
\tag{17}
\]
\underline{[Step 4]}

\subsection*{Step 5 – Plug (17) into the recurrence}
Using (17) in (13) we get the deterministic inequality
\[
f(x_{k+1})-f(x^\star)
\le \bigl(1-\gamma_k\bigr)\bigl[f(x_k)-f(x^\star)\bigr]
   +\frac{L D^{2}}{2}\gamma_k^{2}.
\tag{18}
\]
\underline{[Step 5]}

\subsection*{Step 6 – Choose the canonical step size}
Set $\displaystyle\gamma_k=\frac{2}{k+2}$ as in (5).  
Then $1-\gamma_k = \frac{k}{k+2}$ and $\gamma_k^{2}= \frac{4}{(k+2)^{2}}$.  
Insert into (18):
\[
f(x_{k+1})-f(x^\star)
\le \frac{k}{k+2}\bigl[f(x_k)-f(x^\star)\bigr]
   +\frac{L D^{2}}{2}\cdot\frac{4}{(k+2)^{2}} .
\tag{19}
\]
Simplify the second term:
\[
\frac{L D^{2}}{2}\cdot\frac{4}{(k+2)^{2}}
   =\frac{2L D^{2}}{(k+2)^{2}} .
\tag{20}
\]
Thus
\[
f(x_{k+1})-f(x^\star)
\le \frac{k}{k+2}\bigl[f(x_k)-f(x^\star)\bigr]
   +\frac{2L D^{2}}{(k+2)^{2}} .
\tag{21}
\]
\underline{[Step 6]}

\subsection*{Step 7 – Unfold the recurrence}
Define $h_k:=f(x_k)-f(x^\star)$.  
From (21),
\[
h_{k+1} \le \frac{k}{k+2}h_k + \frac{2L D^{2}}{(k+2)^{2}} .
\tag{22}
\]
Multiplying both sides by $(k+2)(k+1)$ yields
\[
(k+2)(k+1)h_{k+1}
\le k(k+1)h_k + 2L D^{2} .
\tag{23}
\]
Set $A_k:=k(k+1)h_k$.  Then (23) reads
\[
A_{k+1} \le A_k + 2L D^{2}.
\tag{24}
\]
Iterating from $k=0$ to $k=K$ gives
\[
A_{K+1} \le A_0 + 2L D^{2}(K+1).
\tag{25}
\]
Since $x_0\in\mathcal{C}$, by convexity $h_0\le \frac{L D^{2}}{2}$, hence $A_0 =0\cdot1\cdot h_0 =0$.  
Thus $A_{K+1}\le 2L D^{2}(K+1)$.

Recall $A_{K+1} = (K+1)(K+2)h_{K+1}$, so
\[
(K+1)(K+2)\bigl[f(x_{K+1})-f(x^\star)\bigr]
\le 2L D^{2}(K+1) .
\tag{26}
\]
Cancel $(K+1)$ and rearrange:
\[
f(x_{K+1})-f(x^\star) \le \frac{2L D^{2}}{K+2} .
\tag{27}
\]
Rename $k=K+1$ to obtain exactly the bound (6).  
\underline{[Step 7]}

\subsection*{Step 8 – Gap bound}
From (17) we have $\gap_k \le 2\bigl[f(x_k)-f(x^\star)\bigr]$.  
Using (6) yields
\[
\gap_k \le \frac{4L D^{2}}{k+2},
\]
which is (7).  
\underline{[Step 8]}

Thus all statements of Theorem \ref{thm:rate} are proved. $\square$

%%%%%%%%%%%%%%%%%%%%%%%%%%%%%%%%%%%%%%%%%%%%%%%%%%%%%%%%%%%%
\section*{Rate Derivation (With Constants)}
%%%%%%%%%%%%%%%%%%%%%%%%%%%%%%%%%%%%%%%%%%%%%%%%%%%%%%%%%%%%
The constant $2L D^{2}$ in (6) originates from the curvature bound $C_f\le L D^{2}$.
If a tighter curvature $C_f$ is known, the bound improves to
\[
f(x_k)-f(x^\star) \le \frac{2 C_f}{k+2}.
\tag{28}
\]
In the special case where $f$ is $\mu$‑strongly convex, an accelerated variant of Frank–Wolfe can attain a linear rate $O\!\bigl((1-\frac{\mu}{L})^{k}\bigr)$, but the plain method retains the $O(1/k)$ bound derived above.

%%%%%%%%%%%%%%%%%%%%%%%%%%%%%%%%%%%%%%%%%%%%%%%%%%%%%%%%%%%%
\section*{Special Cases}
%%%%%%%%%%%%%%%%%%%%%%%%%%%%%%%%%%%%%%%%%%%%%%%%%%%%%%%%%%%%
\begin{itemize}[leftmargin=*]
\item \textbf{Polytope feasible set.} When $\mathcal{C}$ is a polytope, each $s_k$ is an extreme point; after $k$ iterations $x_k$ is a convex combination of at most $k+1$ vertices (sparsity property).
\item \textbf{Exact line search.} Choosing $\gamma_k = \arg\min_{\gamma\in[0,1]} f\bigl((1-\gamma)x_k+\gamma s_k\bigr)$ yields the same $O(1/k)$ rate, often with a smaller empirical constant.
\item \textbf{Away‑step Frank–Wolfe.} If the algorithm is augmented with away steps, the rate improves to linear under a polyhedral strong convexity condition.
\end{itemize}

\end{document}